% ============================================================
%  Introduzione alla Parte II: Meccanica del corpo rigido
% ============================================================

\chapter*{Introduzione Parte~II}
\addcontentsline{toc}{chapter}{Introduzione alla Parte~II}

La meccanica del corpo rigido è il fondamento su cui poggia l'intera teoria delle strutture.
Prima di introdurre la deformabilità --- che è la caratteristica fisica essenziale di qualsiasi
struttura reale --- occorre comprendere a fondo ciò che il vincolo di rigidità implica: quali
configurazioni sono cinematicamente ammissibili, quante e quali forze sono necessarie per
garantire l'equilibrio, come si caratterizzano le reazioni dei dispositivi che limitano il
movimento. Questo fondamento non è specifico della trave né di alcuna forma geometrica
particolare: vale per corpi di forma qualsiasi, e la sua generalità è una delle ragioni per cui
viene sviluppato prima di qualsiasi specializzazione.

La Parte~I ha fornito gli strumenti: l'algebra delle matrici e i sistemi lineari, la geometria
delle applicazioni lineari e la dualità, i vettori geometrici. La
Parte~II li mette al lavoro. Ogni concetto introdotto come strumento astratto --- i quattro
sottospazi fondamentali, l'operatore aggiunto, l'identità bilineare --- trova qui la sua prima
incarnazione meccanica concreta, e il lettore che ha attraversato la Parte~I con attenzione
riconoscerà, capitolo dopo capitolo, la struttura che era stata anticipata.

\section*{Il programma della parte}

Il \textbf{Capitolo~\ref{cap:cinematica_infinitesima}} affronta la \emph{cinematica
	infinitesima del corpo rigido}. Si introducono i gradi di libertà come proprietà del vincolo
di rigidità, indipendente dall'ipotesi di spostamenti infinitesimi. La formula fondamentale
che lega lo spostamento di un punto generico allo spostamento di un polo e alla rotazione
del corpo è derivata con due approcci indipendenti --- la linearizzazione della rotazione
finita e il principio di invarianza delle distanze --- e presentata nella sua forma matriciale.
Si studiano gli invarianti del campo di spostamento, si introduce l'asse centrale con una
costruzione geometrica e si classifica la varietà dei moti rigidi infinitesimi. La
specializzazione al caso piano conduce alla determinazione del centro di rotazione, che
costituirà il principale strumento geometrico per la classificazione dei sistemi vincolati.

Il \textbf{Capitolo~\ref{cap:statica_CR}} affronta la \emph{statica del corpo rigido} in
perfetto parallelo con il precedente. Punto di partenza sono le equazioni di equilibrio,
che motivano la riduzione di un sistema di forze alle sue grandezze globali: risultante e
momento risultante. Si introduce la coppia di forze, si deriva la legge di variazione del
momento al variare del polo e si studiano gli invarianti del campo dei momenti, con la
determinazione dell'asse centrale del sistema di forze mediante una costruzione geometrica
che rispecchia quella cinematica. La specializzazione al caso piano porta al centro di
riduzione e alla classificazione dei sistemi equivalenti. Il parallelo con il
Capitolo~\ref{cap:cinematica_infinitesima} non è casuale: è la prima manifestazione
visibile della dualità cinematica-statica.

Il \textbf{Capitolo~\ref{cap:dualismo_CR}} porta in primo piano il \emph{dualismo tra
	cinematica e statica}. Le due leggi di distribuzione --- il campo degli spostamenti rigidi
e il campo dei momenti --- hanno la medesima struttura affine, e a ogni concetto cinematico
corrisponde un concetto statico duale. Nel caso piano questo dualismo si manifesta nello
scambio di ruolo tra grandezze scalari e vettoriali che trasforma il centro di rotazione
nell'asse di sollecitazione. Il capitolo introduce il lavoro virtuale come forma bilineare
che accoppia le grandezze cinematiche e statiche, ne calcola l'espressione per forze
concentrate, coppie e forze distribuite, e mostra come la condizione di lavoro virtuale nullo
per ogni spostamento rigido sia esattamente equivalente alle equazioni di equilibrio. Il
principio dei lavori virtuali cessa così di essere un enunciato autonomo e diventa la
traduzione bilineare dell'equilibrio.

Il \textbf{Capitolo~\ref{cap:vincoli_CR}} studia le \emph{prestazioni cinematiche e statiche
	dei vincoli}. Si distingue preliminarmente il vincolo di rigidità --- proprietà interna
del corpo che impone l'uguaglianza di spostamento e rotazione nelle zone di connessione,
con le relative nozioni di distorsione e svincolo --- dai vincoli di collegamento, che sono
l'oggetto principale del capitolo. Per ciascuno dei cinque vincoli elementari --- carrello,
bipendolo, cerniera, glifo e incastro --- si ricavano l'equazione cinematica, il luogo
imposto al centro di rotazione assoluto o relativo e la reazione vincolare, ottenuta
applicando il principio dei lavori virtuali: la reazione è quella grandezza statica che
compie lavoro nullo su tutti gli spostamenti compatibili con il vincolo. Si trattano i
vincoli composti, il cedimento vincolare e i vincoli che collegano più di due corpi.

Il \textbf{Capitolo~\ref{cap:problemi_cin_stat_SCR}} sviluppa la \emph{formulazione
	matriciale dei problemi cinematico e statico} per un sistema di corpi rigidi piani. La
matrice dei vincoli, le cui righe codificano le equazioni cinematiche di ogni vincolo
posizionale, governa entrambi i problemi: il problema cinematico vi appare direttamente,
il problema di equilibrio ne utilizza la trasposta. Questo dualismo --- la matrice
dell'equilibrio è la trasposta della matrice di congruenza --- non è una coincidenza
formale, ma la conseguenza diretta della struttura bilineare del lavoro virtuale.
I quattro sottospazi fondamentali della matrice dei vincoli individuano meccanismi,
stati di autoequilibrio, cedimenti ammissibili e forze non equilibrabili, e la relazione
tra le loro dimensioni generalizza il conteggio dei gradi di libertà del sistema.

Il \textbf{Capitolo~\ref{cap:classificazione_sol_SCR}} ha carattere \emph{operativo}:
fornisce strategie pratiche di soluzione per sistemi di corpi rigidi piani di qualsiasi
tipo, senza costruire esplicitamente la matrice dei vincoli. I problemi cinematico e
statico sono trattati in parallelo, così da rendere visibile il dualismo ad ogni passo
della procedura. Per i sistemi determinati si presentano il metodo grafo-analitico e
il metodo dei corpi liberi; per i sistemi labili e iperstatici si sviluppano le procedure
operative duali --- vincoli fittizi e vincoli ridondanti --- con le relative condizioni
di compatibilità. Si introduce infine la gerarchia strutturale, che consente di
decomporre sistemi complessi in sottosistemi risolvibili in cascata.

Il \textbf{Capitolo~\ref{cap:classificazione_SCR}} sviluppa gli strumenti per
\emph{riconoscere e classificare} i sistemi di corpi rigidi piani. Il conteggio
dei gradi di libertà nominali costituisce un criterio necessario ma non sufficiente:
fissa le dimensioni della matrice di congruenza senza determinarne il rango. A parità
di conteggio, la disposizione geometrica dei vincoli può produrre classificazioni
diverse. Il nucleo del capitolo è un catalogo sistematico delle sottostrutture
elementari più frequenti, da un singolo corpo fino a tre corpi, nel quale i casi
regolari e degeneri sono affiancati e le cause geometriche della degenerazione ---
parallelismo o concorrenza degli assi di vincolo, allineamento dei centri di
rotazione --- sono identificate e interpretate. Il catalogo abilita la classificazione
per decomposizione gerarchica e per riduzione a schema noto, e si raccorda con le
strategie di soluzione del capitolo precedente.

Il \textbf{Capitolo~\ref{cap:TLV}} chiude la parte formalizzando la \emph{triade
	meccanica}: congruenza, equilibrio e principio dei lavori virtuali sono tre enunciati
equivalenti, nel senso che due qualsiasi di essi implicano il terzo. Dal dualismo tra
la matrice di congruenza e la matrice di equilibrio si deriva un'identità bilineare
valida per ogni coppia congruente e ogni coppia equilibrata, che costituisce la forma
matriciale esatta del principio dei lavori virtuali. Dalla triade discendono quattro
corollari operativi: la condizione di compatibilità statica, il principio di
Müller-Breslau per il calcolo delle reazioni, la condizione di compatibilità
cinematica e il metodo del carico unitario per il calcolo degli spostamenti. Il
capitolo ha dunque il carattere di un epilogo sintetico: raccoglie in un'unica
struttura coerente i risultati elaborati nei capitoli precedenti e li eleva a
strumenti generali, pronti per essere trasferiti alla meccanica della trave.

%-------------------------------------------------------------------------------------------------------------------
\section*{Generalità e specializzazione}
%-------------------------------------------------------------------------------------------------------------------

Un aspetto metodologico merita di essere dichiarato esplicitamente. I risultati
sviluppati in questa parte valgono per corpi rigidi di forma qualsiasi: i corpi
sono indicati genericamente come $\mathcal{C}_i$, $\mathcal{C}_j$ senza alcuna
ipotesi sulla geometria. La specializzazione alla trave --- prima come corpo rigido,
poi come corpo deformabile --- avviene nella Parte~III e nel Volume~II rispettivamente.
Questa scelta non è soltanto una questione di ordine espositivo: la generalità della
Parte~II garantisce che la struttura logica della teoria --- dualità, lavoro virtuale,
classificazione dei sistemi vincolati --- sia acquisita una volta per tutte, prima
di qualsiasi ipotesi geometrica. Le specializzazioni successive non aggiungono nuovi
principi: applicano quelli già stabiliti a contesti via via più concreti.

In particolare, il vincolo di rigidità introdotto nel
Capitolo~\ref{cap:vincoli_CR} --- con le nozioni di distorsione e svincolo --- anticipa
già la specializzazione alla trave: le reazioni del vincolo di continuità si
identificheranno con le caratteristiche della sollecitazione, e le distorsioni di
Volterra troveranno in quella struttura la loro collocazione naturale. Questi ponti
tra la meccanica del corpo rigido e la meccanica della trave deformabile saranno
resi espliciti nel momento in cui la specializzazione geometrica lo renderà possibile.

%-------------------------------------------------------------------------------------------------------------------
\section*{Il filo conduttore: la dualità}
%-------------------------------------------------------------------------------------------------------------------

Il filo conduttore di questa parte --- e dell'intera opera --- è la dualità tra
cinematica e statica. La cinematica del corpo rigido descrive il nucleo dell'operatore
di deformazione: gli spostamenti che non producono deformazione interna. La statica
descrive le forze in equilibrio: i sistemi di forze il cui lavoro virtuale è nullo su
tutti gli spostamenti rigidi. Le due descrizioni sono duali nel senso preciso
stabilito dalla Parte~I: a ogni grandezza cinematica corrisponde una grandezza
statica duale, e il loro prodotto ha il significato fisico di lavoro.

Questa simmetria percorre la parte in modo progressivo. Emerge per la prima volta nel
confronto tra i Capitoli~\ref{cap:cinematica_infinitesima} e~\ref{cap:statica_CR},
dove le due trattazioni procedono in parallelo con struttura identica. Viene dichiarata
e formalizzata nel Capitolo~\ref{cap:dualismo_CR}, attraverso il lavoro virtuale come
forma bilineare. Si riflette nella formulazione matriciale del
Capitolo~\ref{cap:problemi_cin_stat_SCR}, dove il dualismo tra la matrice di congruenza
e la matrice di equilibrio è la trascrizione algebrica dello stesso principio. Si
manifesta infine nelle procedure operative del Capitolo~\ref{cap:classificazione_sol_SCR},
dove ogni metodo cinematico ha il suo esatto corrispondente statico. Il
Capitolo~\ref{cap:TLV} raccoglie questo percorso nella triade meccanica: non una
conclusione che chiude, ma una sintesi che apre verso la meccanica della trave.



